% Folder 79, batch 10 (archivists' pages 181-200).
% Pass: Opus 5.5 (claude-opus-5-5), 2026-10-03 — first pass, unchecked against the pages by a human.
%
% The author's own pagination runs on the rectos, top centre: 94 (p. 182)
% to 103 (p. 200). The connective prose is a fast hand and much of it is
% marked doubtful; the formulas carry.
\documentclass[11pt,a4paper]{article}
\input{../preamble/grothendieck}

\folder{79}
\batch{10}
\pages{181}{200}
\foldertitle{2-polyèdres réguliers triangulés, et modules libres de rangs 2 sur les anneaux $\Lambda=\mathbb{Z}/n\mathbb{Z}$ et $\Lambda=\mathbb{Z}$ : notes manuscrites (s.d.).}
\dating{[à partir de 1980]}
\watermark{Édition de démonstration}

\begin{document}

\page{181}
\note{la page reprend au milieu d'un argument commencé avant ce lot : les relations (8), (9), (10) et 13 auxquelles elle renvoie sont sur les pages précédentes.}
deux relations (8) (9), chacune l'une des deux, implique l'autre, i.e.\ l'égalité (10),
i.e.\ le membre \uncertain{de droite} de la relation 13, est en fait une égalité
\[
\text{(15)} \qquad Z_{s_0} = S^{G_{s_0}} .
\]

Voici des cas où il en est ainsi
\[
\text{(16)}\quad
\begin{cases}
\text{a) } S \text{ fini (car tout \uncertain{homomorphisme} surjectif de } S \text{ sur}\\
\qquad \text{lui-même est alors bijectif)}\\
\text{b) les } G_s \text{ finis (raison duale pour les inclusions de la forme (8), (9))}\\
\text{c) } G \text{ est un groupe de Lie réel, et les } G_s \text{ sont fermés}\\
\qquad \text{et n'ont qu'un nombre fini de composantes connexes.}
\end{cases}
\]
\note{les marques d'appareil des lignes a) à c) sont reportées ici : « homomorphisme » est une lecture douteuse ; en c) un mot initial (« les ») est biffé.}

\struck{\ill{}} On peut appeler aussi points \add{$G$-}\emph{antipodiques}
\struck{$s \neq s_0$ tels} d'un $s_0 \in S$ les points $\in S^{G_{s_0}}$, i.e.\ fixés
par tout élément de $G$ qui fixe $s_0$. Dans les cas envisagés a) b) c), ce sont
aussi les points $G$-équivalents à $s_0$, qu'on peut aussi appeler
\emph{strictement antipodiques} à $s_0$. On peut appeler application
\add{$G$-}antipodique dans $S$ toute application qui commute à $G$ (donc
nécessairement surjective) \struck{ainsi} et \add{telle} application
$G$-antipodique une application bijective.

\page{182}
\note{p. 94 de l'auteur.}
Les points antipodiques (resp.\ strict.\ antipodiques) sont donc les images de
$s_0$ par les applications antipodiques (resp.\ strict.\ antipodiques). En un
sens évident : \uncertain{explicitons} ces applications correspond aussi
biunivoquement aux « types » (à $G$-isomorphisme près) de couples $(s_0, s)$ de
points antipodiques (resp.\ strictement antipodiques). On \uncertain{fera}
attention que les relations d'antipodisme entre pts \uncertain{sont} pas
\uncertain{même} une relation d'équivalence, et \uncertain{même} que la relation
stricte l'est — \uncertain{dont} la relation de $C$-équivalence envisagée
\ill{} distinct.

Je m'intéresse maintenant au \uncertain{cas où} $S$ est l'ensemble des sommets
d'un 2-polyèdre \add{connexe} régulier (combinatoire), \struck{et \ill{}
(notée maintenant $G_0$)}, \add{et des s.-g.} groupe des automorphismes, et à la
variante orientée — $G_0$ étant \uncertain{alors} le groupe des automorphismes
orientés. Dans le cas où les ordres des sommets sont pairs,

\page{183}
auquel je vais me limiter maintenant, i.e.\ dans le cas (16) b, il n'y a pas
lieu de distinguer entre \add{la notion d'}antipodisme et celle d'antipodisme
strict.
\marginal{On a \uncertain{vérifié} \uncertain{par} la suite \uncertain{au} cas des polyèdres triangulés (16) b sur les faces des triangles)}

Soit $Z$ le commutant de $G_0$ dans $\mathfrak{S}_S$, \add{i.e.\ le groupe des
antipodismes,} et soit $\lambda \in Z$. On dira que $\lambda$ est une
« homothétie » si on a
\marginal{J'ignore si \uncertain{cette condition est toujours} \uncertain{satisfaite}!}
\[
\text{(17)}\ \text{(H1)} \quad \forall s,t \in S \qquad
\{s,t\} \in A_0 \iff \{\lambda s, \lambda^{-1} t\} \in A_0 ,
\]
où $A_0 \subset \mathfrak{P}_2(S)$ l'ens.\ des arêtes du polyèdre, i.e.\ on
encore, remplaçant \struck{$t$ par $\lambda t$} $s$ par $\lambda^{-1}s$
\[
\text{\struck{i.e.}}\quad
\text{\struck{(18) $\forall s \in S$, \dots\ $\mathrm{Et}_{A_0}($}}
\]
\[
\forall s,t \in S \qquad
\underbrace{\{\lambda^{-1}s, t\} \in A_0}_{\substack{\lambda^{-1}s \in \mathrm{Et}_{A_0}(t),\\ \text{i.e. } s \in \lambda\, \mathrm{Et}_{A_0}(t)}}
\iff
\underbrace{\{s, \lambda^{-1}t\} \in A_0}_{s \in \mathrm{Et}_{A_0}(\lambda^{-1}t)}
\]
i.e.
\[
\text{(18)} \qquad \forall t \in S, \text{ on a } \quad
\lambda\, \mathrm{Et}_{A_0}(t) = \mathrm{Et}_{A_0}(\lambda^{-1} t) .
\]

Ainsi une homothétie transforme l'étoile \uncertain{de} $t$ en celle
\uncertain{(bien définie} — l'étoile de $t$ \uncertain{en tant que centre}
\ill{}) de $\lambda^{-1} t$. La condition (17) \uncertain{équivaut} \ill{} à
l'implication dans le sens $\Rightarrow$
\[
\text{(17$'$)} \qquad \forall s,t \in S \qquad
\{s,t\} \in A_0 \Rightarrow \{\lambda s, \lambda^{-1} t\} \in A_0 ,
\]
et par raison d'homogénéité, comme $\lambda$ commute à $G_0$ qui est transitif
sur l'ens.\ des couples $(s,t)$ tels que $\{s,t\} \in A_0$, il suffit que l'on

\page{184}
\note{p. 95 de l'auteur.}
ait
\[
\text{(17$''$)} \qquad \{\lambda s_0, \lambda^{-1} t_0\} \in A_0
\]
où $(s_0, t_0)$ est une arête orientée fixée. Cela signifie aussi qu'il existe
$g \in G_0$ tel qu'on ait
\[
g(s_0, t_0) = \{\lambda s_0, \lambda^{-1} t_0\}
\]
ou encore, en posant
\[
\text{(19)} \qquad t_0 = \rho(s_0) \qquad \text{pour } \rho \in G_0 \text{ convenable,}
\]
qu'on ait
\[
\text{\struck{$\lambda^{-1} u s_0 = s_0$}} \qquad
g^{-1} \lambda s_0 = s_0 ,
\qquad
(g\rho)^{-1} \underbrace{\lambda^{-1} \rho}_{\rho \lambda^{-1}} s_0 = s_0
\]
\[
\rho^{-1} g^{-1} \rho\, \lambda^{-1} s_0
\]
soit
\[
\text{(20)} \qquad \exists g \in G_0 \text{ tel que } \quad
g^{-1} \lambda s_0 = s_0 , \quad
\rho^{-1}(g^{-1}) \lambda^{-1} s_0 = s_0 .
\]

Explicitons alors $\lambda$ : à l'aide d'un élément
\marginal{Je \uncertain{note maintenant} $\widetilde{U} = G_{0 s_0}$ le \uncertain{stabilisateur} \ill{} ; $U$ : $G^{+}_{0 s_0}$ \ill{}}
\[
\text{(21)} \qquad a \in \mathrm{Norm}_{G_0}(\widetilde{U}) \qquad \lambda s_0 = a s_0 .
\]
La condition (20) s'écrit
\[
\text{(22)} \qquad \exists g \in G_0 \text{ tel que } \quad
g^{-1} a \in \widetilde{U} , \quad \rho^{-1}(g^{-1}) a^{-1} \in \widetilde{U}
\]
où on \struck{peut} prend pour $\rho$ un quelconque élément \uncertain{satisfaisant}
\uncertain{(19)} (déterminé mod $U$ i.e.\ \ill{}) p.\ ex.\ on peut prendre
\add{\uncertain{dans le formulaire} $\mathrm{SL}(2,\mathbb{Z})$ …}
\[
\text{(23)} \qquad \rho = \tau_0 , \ \text{ou } \lambda_1 , \ \text{ou } \lambda_\infty^{-1} = \rho , \ \text{au choix}
\]
\[
\lambda_1 = \tau_0 \tau_\infty , \qquad \lambda_\infty = \tau_0 \tau_1 .
\]
\note{les $\lambda_1$, $\lambda_\infty$ de (23) sont des éléments du formulaire des générateurs (avec $\tau_0, \tau_1, \tau_\infty$), sans rapport avec l'homothétie $\lambda$ ; les deux dernières égalités sont lues de façon douteuse.}

\drawing{en marge gauche, un losange de sommets $u_0$ (en haut), $s_0$, $t_0$, coupé par l'arête $s_0 t_0$ ; dans le triangle supérieur, des marques $\tau_1$, $\tau_0$, $\tau_\infty$ et une petite région hachurée.}

La relation 22 s'écrit aussi, posant $g^{-1} a = u$
\[
\text{(\struck{22})} \qquad \exists u \in \widetilde{U} \text{ tel que } \quad
\rho^{-1}(u a^{-1}) a^{-1} \in \widetilde{U} \qquad \text{i.e.}
\]
\[
a \rho^{-1}(a u^{-1}) \in \widetilde{U}
\qquad \text{i.e.} \quad a \rho^{-1}(a) \in \widetilde{U} \rho^{-1}(u)
\qquad \text{\struck{i.e. $\rho(a) a \in \rho(\widetilde{U}) u$}}
\]
\[
\text{i.e.} \qquad \rho(a) a \in \rho(\widetilde{U}) \dots
\]

\page{185}
ou encore
\[
\rho(a) \text{\struck{$\rho(a)$}}\, a \in \rho(\widetilde{U}) \text{\struck{$\rho^{-1}$}}\, \widetilde{U} .
\]

Posons
\[
\widetilde{U}_0 = \widetilde{U} = G_{0 s_0} , \qquad
\widetilde{U}_1 = G_{0 t_0} = \rho(\widetilde{U}) ,
\]
la condition envisagée s'écrit donc
\[
\text{(24)} \qquad \rho(a) a \in \widetilde{U}_1 \widetilde{U}_0
\]
(en plus de la relation
\[
a \in \mathrm{Norm}_{G_0}(\widetilde{U}_1) \ ).
\]
\note{l'indice 1 de cette dernière relation est lu tel quel ; (21) porte $\mathrm{Norm}_{G_0}(\widetilde{U})$, i.e.\ $\widetilde{U}_0$.}

\emph{NB} Dans le cas où le polyèdre $P_0$ est orienté, il est plus sympathique
dans cette \struck{\ill{}} de travailler dans $G_0^{+}$ (qui est encore
transitif sur l'ens.\ des arêtes orientées, qui est $=$ sur l'ens.\ des \ill{}),
en écrivant \ill{} $=$ un \emph{torseur} sous $G_0^{+}$ \ill{} un élément de
$G_0^{+}$
\[
\text{(25)} \qquad a \in \mathrm{Norm}_{G_0^{+}}(U_0)
\qquad \bigl[\, U_0 = G^{+}_{0 s_0} = \widetilde{U}_0 \cap G_0^{+} \,\bigr]
\]
et en écrivant \struck{la} alors (24) \add{dans $G_0^{+}$,} sous la forme
équivalente
\[
\text{(24$'$)} \qquad \rho(a) a \in U_1 U_0 .
\]

On dit que l'homothétie $\lambda$ est \emph{stricte} si elle satisfait la
condition que la structure de graphe \add{via} $A_\lambda \in \mathfrak{P}_2(S)$,
définie par
\[
\text{(25)} \quad
\begin{cases}
A_\lambda = \bigl\{ \{s, \lambda t\} \in \mathfrak{P}_2(S) \bigm| (s,t) \in S^2 ,\ \{s,t\} \in A_0 \bigr\} \\
\text{i.e. } \ \forall s,t \in S \quad \{s,t\} \in A_0 \iff \{s, \lambda t\} \in A_\lambda
\end{cases}
\]
\note{deux débuts biffés de cette définition (« $\{s,t\} \in$ ») ; le numéro (25) est répété sur la page.}
est isomorphe à $(S, A_0)$, \struck{\ill{}} condition qui

\page{186}
\note{p. 96 de l'auteur.}
\marginal{il faudra y revenir par la suite}
peut aussi s'\add{expliciter} \struck{formuler} en termes de théorie des groupes.
Notons que l'on a alors
\[
\text{(26)} \qquad \mathrm{Et}_{A_0}(s) = \mathrm{Et}_{A_\lambda}(\lambda s)
\qquad \text{pour } s \in S ,
\]
donc les deux structures $A_0$ et $A_\lambda$ donnent lieu aux mêmes étoiles —
mais avec des centres différents : le centre \struck{$s$} de l'étoile $E$, pour
la structure $\Gamma_\lambda = (S, A_\lambda)$, est l'antipodique du centre $s$
pour $\Gamma_0$, via l'antipodisme $\lambda$.
\marginal{cette relation est indépendante de la condition stricte — mais si elle est \uncertain{satisfaite}, la structure $A_\lambda$ \ill{} polyèdre \ill{} (attention : chaque \ill{} $A_\lambda$ \ill{} $\in \mathfrak{P}_2(E)$), \ill{} de polygone \ill{}}

Considérons alors le groupe $G_{0s} = G_{0(s' = \lambda s)}$, \struck{et le sous}
le sous-groupe des rotations relatif \struck{\ill{}} à $s$ : c'est le seul
sous-groupe cyclique d'ordre $n$ de $G_{0s} \simeq \mathbb{D}_n$, donc est égal
au sous-groupe correspondant \struck{pour} relatif à $s'$. Cela montre que sur
l'étoile $s$, les structures de polygone, \uncertain{orientées} induites par
$A_0$ et $A_\lambda$ respectivement, ont même groupe de rotations, i.e.\ se
transforment l'une de l'autre par un \uncertain{élément} de
\[
(\mathbb{Z}/n\mathbb{Z})^{*} / \{\pm 1\} = \mathbb{Z}'_n
\]
(\uncertain{n} \uncertain{étant} l'ordre des points de $S$) ;
\[
\text{(27)} \qquad \alpha_{E,\lambda} = (\alpha_{E,0})^{\nu} , \qquad
\nu \in \mathbb{Z}'_n = (\mathbb{Z}/n\mathbb{Z})^{*} / \pm 1
\]

\page{187}
L'\add{élément} \struck{entier} $\nu$ dépend a priori du choix de $s$, Mais
\struck{si on} par transport de structure, on voit qu'il est le même pour $s$ et
\add{pour} $us$, $u \in G_0$, donc qu'il est en fait indépendant de $s$. Mais
\struck{alors on voit donc que la} la structure $A_\lambda$ est entièrement
déterminée par \struck{l'indice} l'élément
\[
\text{(28)} \qquad \nu \in \mathbb{Z}'_n = (\mathbb{Z}/n\mathbb{Z})^{*} / \pm 1
\]
puisque l'on a
\[
\text{(29)} \qquad A_\lambda = \bigcup_{E \in \Sigma}
\bigl( A_{E,\lambda} = A_{E,0}^{(\nu)} \bigr) .
\]
\note{sous la réunion : « $\Sigma$ ens.\ des étoiles de $S$, \uncertain{les mêmes} pour $A_0$ et $A_\lambda$ ».}

Cet argument montre aussi, en dehors des \uncertain{homothéties} strictes,
\uncertain{en} \uncertain{supposant} la chose suivante \struck{\ill{}}
\[
\text{(30)}\quad
\begin{cases}
\text{b) Toute étoile } E \text{ est, pour la structure } A_{\lambda E} = A_\lambda \cap \mathfrak{P}_2(E),\\
\qquad \text{un polygone combinatoire}\\
\text{a) Toute } a \in A_\lambda \text{ est contenue dans au moins une étoile}
\end{cases}
\]
\note{les lettres b) et a) sont lues dans cet ordre ; une flèche les relie, et la suite les cite comme « a) et b) ».}

Notons que le groupe $G_0$ opère par automorphismes de $A_\lambda$, et qu'il
opère transitivement sur l'ens.\ des arêtes \add{orientées} de $G_0$. Il en
résulte que \add{il suffit de vérifier} les conditions a) et b) pour \emph{une}
étoile $E$, pour \emph{une} arête \struck{\ill{}}. Pour la condition a), il
suffit de prouver dans

\page{188}
\note{p. 97 de l'auteur.}
qu'il existe au moins un « triangle » pour $A_\lambda$, i.e.\ \add{un triplet}
\struck{des sommets $\{s,t,u\}$} mutuellement adjacents pour $A_\lambda$.
\struck{\ill{}} OPS que ce triplet est de la forme \struck{$\{s, \lambda t, u\}$
avec}
\[
\{s, \lambda t, u\} \qquad \text{avec } \{s, \lambda t\} \in A_\lambda
\quad \text{i.e.} \quad \{s,t\} \in A_0
\]
et il reste à vérifier l'existence d'un $u$ qui convient, i.e.\ tel que l'on ait
\[
\{s,u\} \ \& \ \{\lambda t, u\} \in A_\lambda \qquad \text{i.e.}
\]
\[
\{\lambda^{-1} s, u\} , \{t, u\} \in A_0
\]
donc la condition \uncertain{b)} signifie aussi
\[
\text{(31)} \qquad \text{a$'$) si } \{s,t\} \in A_0 , \text{ alors }
\mathrm{Et}_{A_0}(t) \cap \mathrm{Et}_{A_0}(\lambda^{-1} s) \neq \emptyset
\]
condition qu'il suffit de vérifier pour \emph{une} arête $(s,t)$ de sommets
adjacents pour $A_0$. Cela s'énonce aussi \struck{(en \ill{} $s$ et $t$)}
\[
\text{(32)} \qquad \text{a$''$)} \quad \text{\struck{$\forall s \in S$}}
\ \text{(Pour un } s \in S \text{, ou tout } s \in S \text{, c'est kif-kif)}
\]
la distance combinatoire de $\lambda^{-1} s$ à $s$ est $\leq 3$ (et en fait,
égale à 3, \add{\uncertain{si} $\lambda \neq$ \ill{}}, sauf si
$\lambda = $ identité, ou dans le cas de l'octaèdre et $\lambda =$
\struck{antip} l'antipodisme (où elle est 2)).

\drawing{en marge gauche, un chemin de quatre points $s$ — $t$ — $t'$ — $\lambda^{-1}s$ reliés par trois segments.}

Si elle est égale à 0 on a $\lambda^{-1} s = s$ i.e.\ $\lambda = \mathrm{id}$ ;
elle \add{donc 2 sommets adj.} \struck{$s$ et $t$ (adjacents)} ne peut être égale
à 1 puisque \uncertain{stabilisateurs}. Si elle était égale à 2 i.e.\
\struck{$\lambda \neq \mathrm{id}$ et} $\exists t$ adjacent à 2 fois à $s$ et
$\lambda^{-1} s$, i.e.\ on aurait
\[
\mathrm{Et}_{A_0}(s) = \mathrm{Et}_{A_0}(\lambda^{-1} s) \quad \text{et} \quad
s \neq \lambda^{-1} s ,
\]
i.e.\ on est dans le cas de l'octaèdre.

\page{189}
\marginal{On \uncertain{reviendra} \uncertain{donc} \ill{} de la réflexion \ill{} plus \uncertain{tard} …}
Considérons maintenant la condition b), \uncertain{supposant} a) satisfaite.
Si $E = \mathrm{Et}_{A_0}(s)$, alors $E$ \struck{est} stable par
$\widetilde{U}_s$, et l'action de $\widetilde{U}_s$ sur $E$ transforme
$A_{E,\lambda}$ en lui-même.
\note{la page s'arrête ici, laissée blanche au-dessous ; la page suivante recommence au § 2.}

\section{2.}

\page{190}
\note{p. 98 de l'auteur.}
Soit \struck{$(S, A_0)$ un graphe}
\[
\text{(1)} \qquad (S, A_0) \qquad A_0 \subset \mathfrak{P}_2(S)
\]
un graphe, satisfaisant à la condition de « triangularité »
\[
\text{(2)}\quad
\begin{cases}
\text{a) } S = \bigcup_{a \in A_0} a \\
\text{b) } \forall a \in A_0 ,\ a = \{s,t\} ,\ \exists \text{ exactement deux éléments } u \text{ de } S,\\
\qquad \text{tels que } \{s,u\} , \{t,u\} \in A_0
\end{cases}
\]
de telle façon que pour tout $s \in S$, son étoile pour $A_0$
\[
\text{(3)} \qquad \mathrm{Et}_{A_0}(s) = \{ t \in S \mid \{s,t\} \in A_0 \} = E
\]
est de façon canonique munie d'une structure de contour \struck{\ill{}}
\uncertain{induite}, dont l'ens.\ des arêtes est
\[
\text{(4)} \qquad A_{0E} = \mathfrak{P}_2(E) \cap A_0 .
\]
[Je suppose de plus que les contours sont d'ordre $\geq 3$.] Je suppose donné de
plus un groupe d'automorphismes $G_0$ de $(S, A_0)$, transitif \struck{sur} sur
$\vec{A}_0$, et \add{même} \struck{simplement} transitif sur l'ens.\ des repères
triangulaires
\marginal{demandé plus tard}
\[
\text{(5)} \qquad \mathrm{Rep}(S, A_0) \overset{\mathrm{def}}{=}
\bigl\{ (s,t,u) \in S^3 \bigm| \{s,t\} , \{t,u\} , \{u,s\} \in A_0 \bigr\} \ \bigr] .
\]
\struck{Je suppose de plus satisfaite} Considérons l'application
\[
\text{(6)} \qquad \mathrm{Et}_{A_0} : S \longrightarrow \mathfrak{P}(S) , \qquad
s \longmapsto \mathrm{Et}_{A_0}(s)
\]
et soit $\Sigma_A$ son image
\[
\text{(7)} \qquad \Sigma_A = \{ E \in \mathfrak{P}(S) \mid \exists s \in S \text{ t.q. } E = \mathrm{Et}_{A_0}(s) \}
\subset \mathfrak{P}(S) .
\]
On note que pour tout $E \in \Sigma_A$, la structure de contour sur $E$ définie
par (4) ne dépend pas du choix des $s \in S$ tel que $E = \mathrm{Et}_{A_0}(s)$.
On récupère $A_0$, à

\page{191}
partir de ces structures sur les $E \in \Sigma$, par la formule
\[
\text{(8)} \qquad A_0 = \bigcup_{E \in \Sigma} A_{0E} .
\]

Revenant à l'action de $G_0$ sur $S$, $G_0$ opère par transport de structure sur
les \struck{objets} ens.\ introduits, et \struck{les} l'application
\struck{\ill{}} surjection
\[
\text{(9)} \qquad S \longrightarrow \Sigma , \qquad s \longmapsto \mathrm{Et}_{A_0}(s)
\]
commute à l'action de $G$. Les ensembles
\[
S ,\ A_0 ,\ \vec{A}_0 ,\ \mathrm{Rep}(S, A_0) , \ \text{ainsi que}
\]
\[
\text{(10)} \qquad F_0 = \bigl\{ f \in \mathfrak{P}_3(S) \bigm| \forall a \in \mathfrak{P}_2(f) ,\ a \in A_0 \bigr\}
\]
sont des espaces homogènes sous $G$ — tous sauf $\mathrm{Rep}(S,A_0)$ des
espaces homogènes quotients de $\mathrm{Rep}(S, A_0)$, qui est un torseur sous
$G$.

Je m'intéresse maintenant à l'ens.\ $\mathfrak{B}$ $=$ ens.\ des structures de
graphes $(S, A_i)$ sur $S$, \struck{qui} admettant $G$ comme groupe
d'automorphismes, et donnant lieu au \struck{même ensemble}
\add{\uncertain{même ens.}\ d'étoiles} $\Sigma$ que $A_0$ — \struck{Soit}
$\Pi = (A_i)_{i \in \Pi}$ \add{de structures}, je voudrais déterminer cet ens.\
qui a priori ne dépend que de $S$, $G_0 \subset \mathfrak{S}_S$,
\struck{\ill{}} fixé par $G_0$, et de $\Sigma \subset \mathfrak{P}(S)$ —
\struck{\ill{}} la donnée de $A_0 \subset \mathfrak{P}_2(S)$ — cette donnée
assurant uniquement que $\Pi \neq \emptyset$, puisque
\[
\text{(11)} \qquad A_0 \in \Pi \subset \mathfrak{P}(\mathfrak{P}_2(S)) .
\]
Comme on l'a dit, \add{pour $A_0$, un} $A_i$ est déterminé quand on connaît
les structures induites
\[
\text{(12)} \qquad A_{iE} = A_i \cap E \qquad (E \in \Sigma)
\]

\page{192}
\note{p. 99 de l'auteur.}
qui sont des structures de contour sur les $E \in \Sigma$, ces structures
redonnant $A_i$ par
\[
\text{(13)} \qquad A_i = \bigcup_{E \in \Sigma} A_{iE} .
\]
L'application $\mathrm{Et}_{A_i}$ correspondante
\[
\text{(14)} \qquad \mathrm{Et}_{A_i} : S \xrightarrow{\ \text{surj}\ } \Sigma , \qquad
s \longmapsto \mathrm{Et}_{A_i}(s) ,
\]
dépend a priori de \struck{l'indice} la structure choisie $i \in \Pi$.
\struck{Notons que} Comme $G_0$ opère transitivement \add{sur} $\Sigma$ et
respecte la structure $A_i$, \add{les $A_{iE}$} \struck{donc $A_i$} sont connus,
dès qu'on connaît $A_{iE}$ pour \emph{un} $E_0 \in \Sigma$ : l'application
\[
\text{(15)} \qquad \varphi_E : \Pi \longrightarrow \text{structures de contour sur } E_0 ,
\qquad i \longmapsto A_{iE_0} = A_i \cap \mathfrak{P}_2(E_0)
\]
est injective, et son image est formée de structures de contour qui sont
\add{invariées} \struck{stables} par l'opération de $G_{0E_0}$ (stabilisateur de
$E_0$ dans $G_0$) sur $E_0$.

Il est clair d'ailleurs que la \add{donnée d'une} \struck{\ill{} des} structures
de contour sur $E_0$, invariée par $G_{E_0}$, est équivalente à la donnée d'un
système \add{de structures} $(C_E)_E$ de contour sur les $E \in \Sigma$,
invariées par $G$. Je \uncertain{peux} maintenant exprimer, en termes
\struck{de la} \add{d'une} structure de contour $C_{E_0}$ sur $E_0$, invariée
par $G_{0E_0}$, la condition qu'elle provient d'un $i \in \Pi$.

\page{193}
Notons que le fait que $G$ soit transitif sur l'ens.\ des repères de
$(S, A_0)$, signifie aussi \struck{que a) $G$ est transitif sur $S$ (condition
qui ne dépend pas du choix de $A_0$) que pour tout $s \in S$, $G_s$ est}
\[
\text{(17)}\quad
\begin{cases}
\text{a) } G \text{ est transitif sur } S \text{ (condition indépendante du choix de } A_0)\\
\text{b) } \forall s \in S \text{ (ou pour un } s_0 \in S \text{, c'est kif-kif)}\\
\qquad \text{l'opération de } G_s \text{ sur } \mathrm{Et}_{A_0}(s) = E \text{ est transitive}\\
\qquad \text{sur l'ens. des repères de la structure de contour } (E, A_{0E}) .
\end{cases}
\]
\note{dans b), le quantificateur « $\forall$ » est d'abord biffé puis récrit ; « un » est une lecture douteuse.}

Je vais supposer que cette structure de contour, \add{qui est} \add{commune}
\uncertain{à} \emph{tous les cas de $A_0$}, est commune, i.e.\ polygonale,
\marginal{Il signifie donc que \uncertain{cette} structure triangulaire $A_0$ est \uncertain{associée} d'un 2-polyèdre régulier, d'\ill{} \uncertain{quelconque} \ill{} $\mathrm{Et}_{A_0}$ \ill{} …}
de sorte que cette condition implique que
\[
\text{(18)} \qquad D_E = \mathrm{Im}\bigl( G_s \longrightarrow \mathfrak{S}_E \bigr)
\]
est égal au groupe $\mathrm{Aut}(E, A_{0E})$ tout entier, qui est isomorphe au
groupe diédral $\mathbb{D}_n$, et contient un unique sous-groupe
$\Lambda_E = D_E^{+}$ cyclique d'ordre $n$.
\marginal{Je suppose les sommets $s \in S$ d'ordre $n$ pour $A_0$ (\ill{}), \uncertain{que} les \ill{} $E \in \Sigma$ \uncertain{sont} des contours (\uncertain{$n \geq 3$})}

A priori, \add{par sa définition,} ce groupe dépend du choix de $s \in S$,
mais comme $\mathrm{Aut}(E, A_{0E})$ ne dépend \uncertain{de} sa description
comme \ill{} — elle dépend pas \uncertain{du} choix de $s$ — mais de celui de
$A_0$. En fait, on voit qu'il ne dépend pas, mais seulement des données
\marginal{On pourrait \uncertain{faire} \ill{} \uncertain{dès} qu'on \uncertain{vérifie} que $S \to \Sigma$ est \ill{}}
\[
\text{(19)} \qquad \bigl( S ,\ G_0 \subset \mathfrak{S}_S ,\ \Sigma \subset \mathfrak{P}(S) \bigr) ,
\]
\note{dans (19), un mot biffé illisible après « $S$, ».}
car on a aussi
\[
\text{(20)} \qquad D_E = \mathrm{Im}\bigl( G_E \longrightarrow \mathfrak{S}_E \bigr) ,
\]
\uncertain{cela vient de ce qu'}a priori l'image de $G_E \supset G_s$ dans
$\mathfrak{S}_E$ contenant celle de $G_s$, est \add{\uncertain{seulement}}
\struck{\ill{}}

\page{194}
\note{p. 100 de l'auteur.}
pouvait être différente si $S \xrightarrow{\mathrm{Et}_{A_0}} \Sigma$ n'était pas
injective). En effet, l'image de $G_E$ étant à priori comprise entre
$\mathrm{Aut}(E, A_{0E})$ et $\mathrm{Im}(G_s \to \mathfrak{S}_E)$, donc égales
\uncertain{à} celles-ci, puisqu'elles sont égales.

La structure de $n$-polygone \add{$A_{0E_0}$} \struck{sur} $E_0$ est d'ailleurs
entièrement déterminée par la connaissance d'une paire
\[
\text{(21)} \qquad \underline{\omega}_0 = \{ \omega_0 , \omega_0^{-1} \} \subset D_{E_0}
\]
de générateurs de $D_{E_0}$, \struck{opposés} inverses l'un de l'autre,
engendrant $D_{E_0}^{+} = \Lambda_{E_0}$ — c'est donc une paire de permutations
circulaires de $E_0$, dont la connaissance équivaut à celle de la structure
$n$-polygonale $A_{0E_0}$.

\note{tout le passage qui suit, jusqu'au bas de la page, est barré de deux longs traits en croix ; il est transcrit ici, son en-tête « Lemme » encadré.}
\struck{Lemme} Considérons un ensemble $E_0$ de cardinal $n$
($n \in \mathbb{N}$, $n \geq 3$), muni d'une paire (21) de permutations
circulaires définissant une structure $n$-polygonale sur $E_0$,
\struck{pour tout $\nu \in \mathbb{Z}/n\mathbb{Z}$,} \add{\uncertain{de}
\ill{}} la structure \ill{} \add{du torseur sous}
\[
\text{(2\struck{2})} \qquad \Lambda_0 = \Lambda_{\underline{\omega}_0} = \mathbb{Z}/n\mathbb{Z}\, \underline{\omega}_0
\qquad \text{à } \xi \in \Lambda_0 \text{ correspond } T_\xi : E_0 \xrightarrow{\ \sim\ } E_0 .
\]
\struck{différent pour \ill{}} Pour tout $\nu \in \mathbb{Z}/n\mathbb{Z}$
considérons l'homomorphisme de groupes
\[
\text{(23)} \qquad \alpha_\nu : \xi \longmapsto \xi^{\nu} : \Lambda_0 \longrightarrow \Lambda_0 ,
\]
\add{$\Lambda_0$ sur $E_0$, définie par} et l'opération \add{des}
\[
T^{(\nu)}_\xi (x_0) = T_{\alpha_\nu(\xi)}(x_0) = (T_\xi)^{\nu}(x) .
\]

\page{195}
\struck{Considérons} \emph{Proposition} Soit $E$ un torseur sous un groupe
$\Lambda_E \subset \mathfrak{S}_E$ (cyclique d'ordre $n$)
\add{($n \in \mathbb{N}$, \struck{\ill{}} $n \geq 2$)}. Considérons l'ens.\
$\vec{C}(E)$ des structures de contour \add{orienté} \struck{combinatoire} sur
$E$, invariées par $\Lambda_E$. \struck{On a Pour tout Pour tout} Cet ens.\ est
en correspondance biunivoque avec $\Lambda_E$, en associant à tout
$u \in \Lambda_E$ la structure de contour orienté, \uncertain{la seule} pour
laquelle $u$ soit l'opération de rotation \add{génératrice} associée.
\marginal{$\Lambda_E$ est simplement transitif sur l'ens.\ des repères orientés de chacune des structures envisagées (à l'exception des \uncertain{cas} \uncertain{dégénérés} \ill{} 1-gones ou \uncertain{bigones} …)}

En effet, la donnée d'une structure $\vec{C}$ \add{$E$,} \uncertain{orientée} de
contour orienté, sur un ens.\ fini $E$, équivaut à celle d'un
$u \in \mathfrak{S}_E$ — l'opération de rotation \uncertain{génératrice} de
contour orienté. Si \uncertain{un} \ill{} $\vec{C}$ est invariant par
$\Lambda_E$ opérant sur $E$, cela signifie que \add{$u = u^{\ell}$}
\uncertain{l'un}, i.e.\ \ill{}, signifie que $u$ commute \ill{}. Or dans le
cas d'un torseur, le \uncertain{commutant} de \ill{} structure de torseur est
réduit : $\Lambda_E$ — cqfd.
\marginal{\emph{NB} J'ai dû exclure par les contours \uncertain{dégénérés} des orbites d'ordre 1 et 2 — ce qui correspond à une définition \uncertain{inutilement} \uncertain{restrictive} dans l'optique \ill{} la $\Lambda$ \uncertain{structure} est \uncertain{portée} par l'un des $E$ sommets … Si \uncertain{on} \uncertain{voulait} l'éviter, il faudrait dans la proposition \ill{} aux $u \in \Lambda_E$ qui ne sont pas \uncertain{d'ordre} 1 ou 2, ce qui \uncertain{élimine} \ill{} le groupe \ill{} 2 $\Lambda_E$}

\emph{Corollaire 1} \struck{Les} Les structures $\vec{C}$ envisagées forment un
\uncertain{torseur} \uncertain{sous} \ill{} libres de rg 1 sur
$\Lambda_n = \mathbb{Z}/n\mathbb{Z}$ (\uncertain{sous} $\Lambda_E$!). Les
générateurs \ill{} \ill{} structures correspondant aux structures
\uncertain{polygonales} orientées sur $E$, invariées par $\Lambda_E$,
\add{qui forment donc un} \emph{torseur} \uncertain{sous}
$\Lambda_E^{*}$ \uncertain{$/$sous} $(\mathbb{Z}/n\mathbb{Z})^{*}$.

\page{196}
\note{p. 101 de l'auteur.}
\emph{Corollaire 2} Toute structure \struck{orientée} de contour \add{$C$} sur
$E$, invariée par $\Lambda_E$, admet exactement deux orientations invariées par
$\Lambda_E$, elles sont \uncertain{opposées} l'une à l'autre. De cette façon,
l'ensemble $\mathcal{C}(E)$ de ces contours s'identifie à l'ensemble
\add{\uncertain{impair}}
\[
\mathcal{C}(E) \simeq \bigl( \Lambda_E \smallsetminus \{0, \sigma\} \bigr) / \pm \mathrm{id}_{\Lambda_E}
\]
(\struck{\ill{}} \uncertain{dans} l'ensemble $\Lambda_E / \pm \mathrm{id}_{\Lambda_E}$
plus sympathique, \uncertain{il faut} \uncertain{ajouter} les deux contours
\uncertain{invariés} « \uncertain{impropres} », correspondant à $0$ et
$\sigma$ …) dans ce cas, \uncertain{le} \uncertain{monoïde}
$\Lambda_n^{\times} = \mathbb{Z}/n\mathbb{Z}$ avec la multiplication
\struck{\ill{}} \struck{le groupe} opère sur cet ensemble, qui \add{est}
\struck{admet un} \add{\ill{}} \struck{\ill{}} \uncertain{non com.}$^{t}$ ;
$\Lambda_n^{\times}$ \uncertain{monoïde} \uncertain{opérateurs}
\struck{\ill{}}.
\marginal{\emph{NB} Le \uncertain{sous-groupe} \ill{} de $\mathfrak{S}_E$ qui \uncertain{normalise} \ill{} $\Lambda_E$ est \ill{}, i.e.\ \ill{} $\mathcal{C}(E)$ \ill{} $\pm 1 \subset$ \ill{} (22) … \ill{} $\Lambda_E$ \ill{} $\mathrm{Aff}(\Lambda_E)$ \ill{} $D_E$ \ill{}}

L'ens.\ des structures de polygones sur $E$ \uncertain{invariées} par $E$
\struck{est un espace 1-1 ! \ill{} un torseur} \add{forme un} torseur sous
\struck{\ill{}} $\Lambda_E^{*} / \pm 1$, et forme
\[
\text{(23)} \qquad \mathbb{Z}'_n = (\mathbb{Z}/n\mathbb{Z})^{*} / \pm 1 .
\]
Pour $\nu \in (\mathbb{Z}/n\mathbb{Z})^{*}$ resp.\
$\dot{\nu} \in \mathbb{Z}'_n = (\mathbb{Z}/n\mathbb{Z})^{*} / \pm 1$,
\uncertain{désignons} par $T_\nu$ resp.\ $T_{\dot{\nu}}$ l'opération
\uncertain{induite} par $\nu$ resp.\ sur $\vec{C}_E$ resp.\ $\hat{C}_E$ —
correspondant sur l'un à la \ill{}, et cette \uncertain{opération} \ill{} les
\uncertain{éléments} $\nu \in (\mathbb{Z}/n\mathbb{Z})^{*}$ (\uncertain{dans} $*$!)
\uncertain{peut} être prolongée au cas des
$\dot{\nu} \in (\mathbb{Z}/n\mathbb{Z}) / \pm 1$ …

\page{197}
En termes d'un générateur $\omega_0$ de $\Lambda_E$, i.e.\
$\omega_0 \in \Lambda_E^{*}$, on trouve la structure de contour orienté
correspondant à $\nu \cdot [\omega_0] = [\omega_0^{\nu}]$, i.e.\
$u = \omega_0^{\nu}$ — et pour $\dot{\nu} = \{\pm \nu\}$, on trouve la structure
de contour définie par la paire $\{ \omega_0^{\nu} , \omega_0^{-\nu} \}$.

Revenons maintenant au cas de la situation
\[
\bigl( S ,\ G_0 \subset \mathfrak{S}_S ,\ \Sigma \subset \mathfrak{P}_E \bigr) ,
\]
\marginal{\emph{NB} On \uncertain{voit} maintenant \uncertain{que} $G$ \uncertain{est} transitif \uncertain{sur} l'ens.\ des \ill{}}
\uncertain{on} voit donc que les systèmes
\[
(C_E)_{E \in \Sigma} \qquad
C_E \text{ structure de contour sur } E ,\ C_E = (E, A_E)
\]
invariés par $G$, correspondent biunivoquement aux éléments de
\[
\bigl( \Lambda_{E_0} \smallsetminus \underbrace{\{0, \sigma\}}_{\text{sous-groupes d'ordre 2}} \bigr) / \pm 1 .
\]
\struck{Nous} Nous proposons maintenant, \struck{pour} un tel $\nu_{E_0}$ ou
$\nu_E$, d'examiner s'il provient d'une \add{structure} $i \in \Pi$.
\marginal{L'ens.\ \uncertain{symétrique} \ill{} d'éléments $(\Lambda_E \smallsetminus \{0, \sigma_E\}) / \pm 1$ \ill{} $\in (\Lambda_E$ \ill{} …}

\struck{Supposons donc} Considérons donc le système \struck{$(C_E)$}
$(\nu_E)_{E \in \Sigma}$ lié à un $\nu_{E_0}$ associé à $C_E = (E, A_E)$ et
posons
\[
\text{(24)} \qquad A = \bigcup_{E \in \Sigma} A_E \subset \mathfrak{P}_2(S) ,
\]
il est clair que $A$ est stable sous $G$, \add{et que} $G$ est transitif sur
l'ens.\ des arêtes orientées de la structure de graphe $(S, A)$.

\section{3.}

\page{198}
\note{p. 102 de l'auteur.}
Je vais reprendre la notion \add{de} graphe et de graphe triangulaire, dans une
optique un peu différente. Une structure de graphe strict peut être interprétée
comme la donnée d'un ens.\ de sommets $S$ et d'une partie \emph{symétrique} de
$S \times S \smallsetminus S$
\[
\text{(1)} \qquad R_0 \subset S \times S , \qquad R_0 \cap \mathrm{diag} = \emptyset , \qquad {}^{s}R = R ,
\]
l'étoile d'un $s \in S$ s'écrivant alors comme
\[
\text{(2)} \qquad \mathrm{Et}_{A_0}(s) = R_0(s) \qquad \text{on note aussi } \mathrm{Et}_{R_0}(s) \ ?
\]
\marginal{\emph{NB} si on n'exige pas $R_0 \cap \mathrm{diag} = \emptyset$, on a une notion un peu plus générale, avec \ill{} $(S)$, \ill{} $A_0 \subset \mathfrak{P}$ \ill{} conditions \ill{} … des \uncertain{boucles} \ill{}}
pour l'ens.
\[
\text{(3)} \qquad A_0 = \bigl\{ \{s,t\} \in \mathfrak{P}_2(S) \bigm| (s,t) \in R_0 \bigr\} .
\]
\struck{Dans le cas où} Considérons l'application
\[
\text{(4)} \qquad S \longrightarrow \mathfrak{P}(S) , \qquad s \longmapsto R_0(s) = \mathrm{Et}_{A_0}(s)
\]
\struck{et injection,} dirigeons \struck{\ill{}} par $S^{*}$ son
\add{image} \struck{\ill{}} (notée $\Sigma$ précédemment)
\[
\text{(5)} \qquad S^{*} = \{ E \in \mathfrak{P}(S) \mid \exists s \in S ,\ \text{t.q. } E = R_0(s) \} ,
\]
on note que l'inclusion identique
\[
\text{(6)} \qquad S^{*} \hookrightarrow \mathfrak{P}(S) , \qquad
s^{*} \longmapsto E_{s^{*}}
\]
\uncertain{définissant} une relation s'interprète aussi comme définissant une
relation
\[
\text{(7)} \qquad \vec{R}_{*} \subset S \times S^{*}
\]
lire :
\[
\text{(8)} \qquad \forall s^{*} \in S^{*} , \text{ on a } \quad {}^{s}R_{*}(s^{*}) = E_{s^{*}} .
\]

[\struck{d'où} Nous nous intéressons surtout au cas où
\struck{(9) $\forall s \in S$, on a $R_0(s)$} (\add{$A_0$-}étoile détermine son
centre), (4) est injection \struck{\ill{}} alors, prenant tout
$s \in S^{*}$, un \ill{} \struck{\ill{}} application bijection canonique
\[
\text{(9)} \qquad \underbrace{R_{*}(s)}_{\substack{\text{ens. des étoiles}\\ E \text{ qui contiennent } s}}
\xrightarrow{\ \sim\ } \mathrm{Et}_{A_0}(s)
\]
\[
t \longmapsto \mathrm{Et}_{A_0}(t)
\]
\note{sous la flèche (9), une flèche qui associe à $t$ l'étoile $\mathrm{Et}_{A_0}(t)$, tracée de droite à gauche.}
(en général, on a seulement une application surjective)]

\page{199}
\struck{Posons} Écrivons maintenant \uncertain{indiciellement} l'application (4)
:
\[
\text{(10)} \qquad \mathrm{Et}_{A_0}(s) = E_{\varphi_{A_0}(s)} = {}^{s}R_{*}(\varphi(s))
\]
où
\[
\text{(11)} \qquad \varphi_{A_0} : S \longrightarrow S^{*}
\]
est l'application surjective [en fait bijection \add{sous l'hypothèse
envisagée plus haut} \struck{faite}) \ill{} induite par (4),
\struck{\ill{}}
[\struck{\ill{}} Ainsi la formule (9) s'écrit aussi
\[
\text{(12)} \qquad \forall s \in S , \qquad
\underset{\substack{\cap \\ S^{*}}}{R_{*}(s)} \xrightarrow{\ \sim\ } E_{\varphi_0(s)} .\ ]
\]

Ainsi, la structure \add{de graphe strict} $(S, A)$,
$A \subset \mathfrak{P}_2(S)$, [satisfaisant la condition d'injectivité
\[
\text{(13)}\quad
\begin{cases}
s \longmapsto \mathrm{Et}_{A_0}(s) \text{ de } S \text{ dans } \mathfrak{P}(S) \text{ injective}\\
\text{i.e. } \forall s,t \in S ,\ s \neq t ,\ \exists u \in S \text{ tel que l'on ait}\\
\qquad (\{s,u\} \in A_0 ,\ \{t,u\} \notin A_0) \text{ ou l'inverse}\\
\qquad (\{t,u\} \in A_0 ,\ \{s,u\} \notin A_0)\\
\text{i.e. } \forall a \in \mathfrak{P}_2(S) ,\ \exists u \in S ,\ s,t \in a ,\ \text{tels que}\\
\qquad \{s,u\} \in A_0 ,\ \{t,u\} \notin A_0 \ ]
\end{cases}
\]
définit les structures suivantes
\marginal{\emph{NB} Même sans condition (13), la structure $(S, A_0)$ définit (14) \ill{} ; cela n'est plus \uncertain{vrai} que $\varphi$ est bijection, \uncertain{elle est} seulement surjective — \uncertain{mais} \ill{} …}
\[
\text{(14)}\quad
\begin{cases}
\text{a) une partie } S^{*} \subset \mathfrak{P}(S) , \text{ donc un ens. } S^{*}\\
\qquad \text{et une relation } R_{*} \subset S \times S^{*}\\
\text{b) une application [bijection] } \varphi_0 : S \longrightarrow S^{*} ,
\end{cases}
\]
\add{la relation $R_0 \subset S \times S$,} permettant de reconstruire
\add{donc la} structure de graphe $(S, A_0)$, par la formule
\marginal{Mais \ill{} la condition \ill{} …}
\[
\text{(15)} \qquad \forall s,t \in S , \qquad
(s,t) \in R_0 \iff (s, \varphi_0(t)) \in R_{*}
\]
\[
\bigl[\text{i.e. } s \in \mathrm{Et}_{A_0}(t) = {}^{s}R_0(\varphi(t)) \bigr]
\]
i.e.
\[
\text{(15$'$)} \qquad R_0 = (\mathrm{id}_S \times \varphi_0)^{-1}(R_{*}) ,
\]
d'où
\[
\text{(16)} \qquad A_0 = \bigl\{ a \in \mathfrak{P}_2(S) \bigm| \exists (s,t) \in S ,
\text{ avec } a = \{s,t\} ,\ (s, \varphi_0(t)) \in R_{*} \bigr\} .
\]

\page{200}
\note{p. 103 de l'auteur.}
Pour que la structure
\[
\text{(17)} \qquad (S, S^{*}, R_{*}, \varphi_0)
\]
\[
\text{(18)}\quad
\begin{array}{l}
S \text{ ens. des « sommets »}\\
S^{*} \text{ ens. des « étoiles » ou « cosommets »}\\
R_{*} \text{ relation « d'incidence »}\\
\varphi_0 \text{ application jauge}
\end{array}
\qquad A_0 \subset \mathfrak{P}_2(S)
\]
\note{dans la liste (18), « application » est d'abord écrit puis biffé avant « application jauge ».}
provienne, à isom.\ près, d'une structure de graphes sur $S$,
\struck{i.e.\ d'une} satisfaisant (13), il faut évidemment qu'elle satisfasse
les conditions suivantes
\marginal{\emph{NB} \uncertain{Si} \uncertain{on} \uncertain{exige} que $(S, A_0)$ satisfasse la condition (13), il \uncertain{faut} laisser l'injectivité dans a) et dans b) \ill{} $S \to \mathfrak{P}(S^{*})$, \ill{} de $\varphi_0$ \uncertain{seulement} surjective.}
\[
\text{(19)}\quad
\begin{cases}
\text{a) les applications}\\
\qquad s^{*} \longmapsto {}^{s}R_{*}(s^{*}) : S^{*} \longrightarrow \mathfrak{P}(S)\\
\qquad s \longmapsto R_{*}(s) : S \longrightarrow \mathfrak{P}(S^{*})\\
\qquad \text{sont injectives}\\
\text{b) L'application } \varphi_0 : S \longrightarrow S^{*} \text{ est bijection, et}\\
\qquad \text{satisfait la condition de symétrie :}\\
\qquad \text{la relation } (s, \varphi_0(t)) \in R_{*} \text{ en } (s,t) \text{ est symétrique}\\
\qquad \text{i.e. } \forall s,t \in S , \ \text{équivalence}\\
\qquad (s, \varphi_0(t)) \in R_{*} \iff (t, \varphi_0(s)) \in R_{*} .\\
\text{c) } \forall s \in S , \ (s, \varphi_0(s)) \notin R
\end{cases}
\]
\note{en c) un début « la relation envisagée » est biffé.}
\marginal{\emph{NB} On peut \struck{expliciter} la donnée de $\varphi_0$, plus symétriquement, par la donnée d'une 2e relation $\Phi_0 \subset S \times S^{*}$, bijective …}

(L'injectivité de $\varphi_0 : S \to S^{*}$ i.e.\ de l'application composée
$S \to S^{*} \to \mathfrak{P}(S)$ résulte de l'hypothèse d'injectivité (13).)
La réciproque est évidente. \struck{Nous avons distinct} Ainsi, on trouve

\[
\text{(20)}\quad \text{Équivalence de groupoïdes}
\]
\[
\begin{array}{c}
\text{(A)}\\
\text{catégorie-groupoïde}\\
\text{des graphes stricts}\\
(S, A_0) \text{ satisfaisant (13)}\\
\text{(sommets déterminés}\\
\text{par leurs étoiles)}
\end{array}
\ \xrightarrow{\ \approx\ }\
\begin{array}{c}
\text{(B)}\\
\text{catégorie-groupoïde des}\\
\text{systèmes } (S, S^{*}, R_{*})\\
\text{« sommets-cosommets-incidence »}\\
\text{satisfaisant les conditions d'inj. symétriques,}\\
\text{munis d'une application bijection}\\
S \xrightarrow{\varphi_0} S^{*} \text{ « symétrique »}\\
\text{et « antiréflexive » } (s, \varphi_0(s)) \notin R_0
\end{array}
\]

La description (B) fait apparaître une \emph{dualité} dans la catégorie des
graphes envisagés (« graphes stricts séparés »), via
\marginal{« séparé » signifie que \uncertain{la notion} d'étoile \uncertain{permet} \uncertain{de retrouver} d'un sommet \ill{} de ses sommets}
\[
\text{(21)} \qquad (S, S^{*}, R_{*}, \varphi_0) \longmapsto
\bigl( \underset{S'}{S^{*}} ,\ \underset{S'^{*}}{S} ,\ \underset{R'_{*}}{{}^{s}R_{*}} ,\ \underset{\varphi'_0}{\varphi_0^{-1}} \bigr) ,
\]
d'où la notion de « graphe dual » d'un graphe strict séparé mais celui-ci est
canoniquement isomorphe \add{au graphe de départ}, \uncertain{via}
\[
\text{(22)} \qquad S \xrightarrow{\ \varphi_0\ } S^{*} \qquad S^{*} \xrightarrow{\ \varphi_0^{-1}\ } S
\]
\marginal{\emph{NB} La condition de symétrie (19, b) signifie justement que (22) \uncertain{est} \uncertain{compatible} \uncertain{avec} les relations …}
\note{l'argument se poursuit au-delà de ce lot.}

\end{document}
